% Susunan definisi cloud-native CNCF v1.1 dibaca dari kiri ke kanan: praktik,
% contoh teknologi, sifat sistem yang dituju, lalu hasil yang diharapkan.
\begin{tikzpicture}[
  font=\scriptsize,
  execute at begin node={\hyphenpenalty=10000},
  kotak/.style={draw, rounded corners, align=left, text width=2.25cm,
                minimum height=3.3cm, anchor=north, inner sep=4pt},
  judul/.style={font=\scriptsize\bfseries, align=center, text width=2.5cm,
                anchor=south},
  panah/.style={-{Stealth[length=2mm]}, thick}
]
  \node[kotak, fill=praditagreen!15] (praktik) at (0,0)
    {mengembangkan, membangun, dan men-\textit{deploy} \textit{workload}\\[3pt]
     di \textit{cloud} publik, privat, atau hibrida\\[3pt]
     pada skala besar\\[3pt]
     terprogram dan dapat diulang};
  \node[kotak, fill=blue!8] (teknologi) at (3,0)
    {\textit{container}\\ \textit{service mesh}\\ \textit{multi-tenancy}\\
     \textit{microservices}\\ \textit{immutable infrastructure}\\
     \textit{serverless}\\ API deklaratif\\[3pt]
     {\tiny\mdseries daftar tidak lengkap}};
  \node[kotak, fill=orange!12] (sifat) at (6,0)
    {keterkaitan longgar\\[2pt] aman\\[2pt] tangguh\\[2pt] terkelola\\[2pt]
     berkelanjutan\\[2pt] dapat diamati};
  \node[kotak, fill=green!10] (hasil) at (9,0)
    {perubahan berdampak besar\\[3pt]
     dilakukan sering\\[3pt]
     hasilnya dapat diprediksi\\[3pt]
     kerja manual berulang minimal};

  \node[judul] at (praktik.north) {Praktik};
  \node[judul] at (teknologi.north) {Teknologi (contoh)};
  \node[judul] at (sifat.north) {Sifat sistem};
  \node[judul] at (hasil.north) {Hasil};

  \draw[panah] (praktik.east) -- (teknologi.west);
  \draw[panah] (teknologi.east) -- (sifat.west);
  \draw[panah] (sifat.east) -- node[above, font=\tiny, align=center]
    {+ otomatisasi} (hasil.west);
\end{tikzpicture}
